\chapter{Practical Filtering}

In the last chapter we designed filters: we chose a cutoff, picked a window or a
Butterworth shape, and ended up holding a set of coefficients. That is the easy
part. A filter on paper is an idealization --- an infinitely precise impulse
response, applied to an infinitely long signal, with infinite numerical
accuracy, in zero time. The moment you want to run it on a stream of samples
arriving in real time from an accelerometer strapped to a moving wrist, every
one of those idealizations breaks, and a new set of questions appears that the
textbook filter-design chapter never asked:

\begin{itemize}
  \item The signal never ends. How do I filter \emph{now}, with only the
        samples I have seen so far, before the next one arrives?
  \item Samples come one at a time, but the FFT wants a whole block. Do I
        filter sample-by-sample, or buffer up a block and transform it? What is
        the cost of each?
  \item A long FIR filter convolved the slow way costs $N$ multiplies per output
        sample. Can the FFT make that cheaper, even on a live stream?
  \item My multiplies are not exact --- they round. On a microcontroller they
        may not even be floating point. When does that rounding bite?
  \item Every filter delays the signal. For a musical instrument you are
        \emph{playing}, that delay is felt in the fingertips. How much can I
        afford, and where does it come from?
  \item Real sensor data has spikes, dropouts, and drift that no linear filter
        handles gracefully. What actually cleans it up?
\end{itemize}

This chapter is about turning a designed filter into a thing that runs. It is
the most concrete chapter in the book, and one of the most important for the
wrist controller, because a gesture instrument lives or dies on \emph{latency}
and on \emph{robust de-noising}. We will build everything from the same
geometric pictures we have been using --- convolution as a sliding overlap, the
spectrum as a set of frequency knobs --- and then get our hands dirty with
buffers, ring indices, fixed-point scaling, and the humble moving average that
does ninety percent of the real work in a sensor pipeline.

\section{Streaming: the signal that never ends}
\label{ch17:streaming}

When we wrote convolution in Chapter~13 we treated the whole signal as a vector
sitting in memory. That is \emph{offline} or \emph{batch} processing: you have a
recording, you process it, you save the result. Plenty of DSP is like that ---
cleaning up a finished audio file, analyzing a stored IMU log. But a live
instrument is the opposite. Samples arrive forever, one (or a few) at a time,
and you must emit output continuously, keeping up with the input rate and never
falling behind.

\begin{definition}[Real-time / streaming processing]
A processing system is \emph{real-time} if it produces each output within a
fixed, bounded time of receiving the input that determines it, and \emph{keeps
up} indefinitely: the average time to process one sample must be less than the
sample period $T_s = 1/f_s$. A \emph{streaming} algorithm is one that can run on
an unbounded input using bounded memory, holding only a finite window of past
samples (its \emph{state}).
\end{definition}

The two words are not the same. ``Streaming'' is about memory: you cannot store
the whole signal, so you keep only what you need. ``Real-time'' is about a
deadline: you must finish before the next sample, on average, with a worst case
you can bound. A filter can be streaming but too slow for real time (it keeps
up on a fast PC but not on a tiny microcontroller), and offline processing can
be wonderfully fast without being real-time at all.

\begin{intuitionbox}
Picture the signal as a paper tape being pulled past a reading head at $f_s$
characters per second. Offline processing is reading the whole tape spread out
on a table. Streaming is reading it through a slot: you only ever see a short
window, and once a character scrolls past the slot it is gone unless you copied
it. The filter's memory is exactly the part of the tape currently inside the
slot --- and a good filter design keeps that slot short.
\end{intuitionbox}

\subsection{What a causal filter is allowed to know}

A filter that runs on a stream can only use samples that have already arrived.
That is the \emph{causality} constraint from Chapter~12, and now it is not a
mathematical nicety --- it is physics. At sample time $n$ you may use
$x[n], x[n-1], x[n-2],\dots$ and past outputs $y[n-1],y[n-2],\dots$, but you
cannot use $x[n+1]$, because it does not exist yet.

This immediately rules out the symmetric, zero-delay filters we like to draw.
The ideal low-pass filter has a $\sinc$ impulse response that stretches
infinitely in \emph{both} directions; half of it is in the future. A symmetric
linear-phase FIR of length $N$ is centered at tap $(N-1)/2$, so to compute the
output ``at the center'' you need samples up to $(N-1)/2$ ahead. In a stream you
cannot peek ahead, so you simply wait: you emit the centered result $(N-1)/2$
samples \emph{late}. That delay is the price of causality, and we will track it
obsessively in Section~\ref{ch17:latency}.

\begin{focusbox}
Three knobs control everything in this chapter, and they trade off against each
other:
\begin{itemize}
  \item \textbf{Quality} --- how sharp the cutoff, how deep the stopband.
  \item \textbf{Latency} --- how long after a sample arrives you can act on it.
  \item \textbf{Cost} --- multiplies, adds, and memory per sample.
\end{itemize}
You cannot maximize all three. A sharp linear-phase FIR is high quality but
high latency; a short IIR is low latency and cheap but has phase distortion. The
art of practical filtering is picking the right corner of this triangle for the
job.
\end{focusbox}

\section{Sample-by-sample versus block processing}
\label{ch17:sample-vs-block}

There are two fundamental ways to structure a streaming filter, and the choice
colors the whole implementation.

\paragraph{Sample-by-sample.} You process exactly one input sample at a time,
producing one output sample, updating a little bundle of internal state in
between. This is how you naturally implement a difference equation, an IIR
filter, a one-pole smoother, a PLL. The code reads like the math:

\begin{lstlisting}[language=C]
// Direct-Form I FIR, one sample in -> one sample out.
// h[0..N-1] are the taps; z[] is a ring buffer of the last N inputs.
float fir_step(float x, const float *h, float *z, int *pos, int N) {
    z[*pos] = x;                 // store newest input
    float acc = 0.0f;
    int k = *pos;
    for (int i = 0; i < N; i++) {
        acc += h[i] * z[k];      // newest tap times newest sample, etc.
        k = (k == 0) ? N - 1 : k - 1;   // walk backward in time
    }
    *pos = (*pos + 1) % N;       // advance ring head
    return acc;
}
\end{lstlisting}

The virtues are minimal latency (you can emit output the instant a sample
arrives) and simple, predictable timing. The cost is that long FIR filters do
$N$ multiply-accumulates (MACs) per sample, every sample, forever.

\paragraph{Block (frame) processing.} You collect $L$ samples into a buffer,
then process the whole block at once. This is mandatory for anything built on
the FFT --- spectrograms (Chapter~10), fast convolution (below), spectral
subtraction --- because the FFT operates on a block. It is also how audio
hardware actually delivers data: a sound card or an I2S peripheral hands you a
buffer of (say) 64 or 128 samples per interrupt, not one sample per interrupt,
because per-sample interrupts would drown the CPU.

\begin{happeningbox}
Block processing buys efficiency by amortizing fixed costs --- the FFT's
$O(L\log L)$, function-call overhead, cache warm-up --- across many samples. But
it spends \emph{latency} to do it: you cannot start processing a block until its
last sample has arrived, so the first sample of a block waits one full block
time, $L\,T_s$, before anything happens. Bigger blocks are more efficient and
more laggy. This single tradeoff is why a guitar amp plugin set to a 512-sample
buffer at 48~kHz feels rubbery (about 11~ms of buffering latency before any
processing) while 32 samples feels immediate but burns CPU.
\end{happeningbox}

\begin{table}[H]
\centering
\begin{tabular}{@{}lll@{}}
\toprule
 & \textbf{Sample-by-sample} & \textbf{Block} \\
\midrule
Latency        & minimal (filter delay only) & block delay $+$ filter delay \\
FFT-based work & not possible                & natural \\
Per-sample cost& fixed, every sample         & amortized, bursty \\
Code shape     & difference equation         & buffer then transform \\
Typical use    & IIR, smoothers, control     & fast convolution, spectral \\
\bottomrule
\end{tabular}
\caption{The two streaming structures. Most real systems mix them: a
sample-by-sample pre-filter feeding a block-based analysis stage.}
\label{ch17:tab-sampleblock}
\end{table}

\begin{goalbox}
For the wrist controller, the IMU runs slowly --- typically $f_s$ between
$100$ and $1000$~Hz --- so per-sample MACs are cheap and we will favor
\emph{sample-by-sample} processing on the motion path to keep latency near
zero. The EMG/FMG path, if present, samples faster (1--2~kHz) and may want short
blocks for feature extraction. The lesson to carry forward: keep the
\emph{control} path (the part that decides ``a note happened now'')
sample-by-sample and short; push any heavy spectral analysis onto a
block-based path that is allowed to lag, and use it only for slow decisions like
``which gesture mode are we in.''
\end{goalbox}

\section{Fast convolution: filtering with the FFT}
\label{ch17:fastconv}

Suppose your filter is genuinely long --- an FIR with $N = 1024$ taps, say a
finely-shaped equalizer or a reverberation impulse response. Direct convolution
costs $N$ MACs per output sample. Can we do better? Yes, and the lever is the
convolution theorem from Chapter~11: convolution in time is multiplication in
frequency.

\begin{equation}
  y[n] = (x \conv h)[n]
  \quad\Longleftrightarrow\quad
  Y(e^{\jj\omega}) = X(e^{\jj\omega})\,H(e^{\jj\omega}).
  \label{ch17:convthm}
\end{equation}

So we could FFT the signal, multiply by the FFT of $h$, and inverse-FFT. The cost
of an FFT of length $M$ is $O(M\log M)$, so multiplying two spectra and
transforming back is far cheaper than $N$ multiplies per sample when $N$ is
large. There is just one catch, and it is a sharp one.

\subsection{The wrap-around problem}

The DFT does not implement ordinary (linear) convolution. It implements
\emph{circular} convolution: it treats both sequences as if they were periodic,
wrapped around a circle of length $M$. When you multiply two length-$M$ DFTs and
invert, the output is

\begin{equation}
  y_{\text{circ}}[n] = \sum_{k=0}^{M-1} x[k]\, h[\,(n-k)\bmod M\,].
  \label{ch17:circ}
\end{equation}

The $\bmod M$ is the villain. The tail of $h$ that should have spilled past the
end of the block instead wraps around and corrupts the \emph{beginning}, a defect
called \emph{time-domain aliasing}. The fix is classic: give the convolution
enough room so the tail has somewhere to go that is not the start.

\begin{whybox}
Linear convolution of a length-$L$ signal block with a length-$N$ filter
produces an output of length $L + N - 1$. If we run the DFT at a length
$M \ge L + N - 1$ --- by zero-padding both inputs out to $M$ --- then the
wrapped copies of the result land entirely in the zero-padding region and never
overlap the genuine output. Circular convolution at length $M \ge L+N-1$
\emph{equals} linear convolution. That single inequality is the whole secret of
fast convolution.
\end{whybox}

\subsection{Overlap-add}
\label{ch17:ola}

A stream is too long to transform in one shot, so we chop it into consecutive
blocks of length $L$ and filter each block separately. Each block, convolved
with the length-$N$ filter, grows to length $L+N-1$ --- so consecutive outputs
\emph{overlap} by $N-1$ samples. Overlap-add says: just add the overlapping
tails back together.

\begin{equation}
  x[n] = \sum_{r} x_r[n], \qquad
  y[n] = \sum_r (x_r \conv h)[n],
  \label{ch17:olasum}
\end{equation}
where $x_r$ is the $r$-th block of $L$ input samples (zeros elsewhere). Because
convolution is linear, the convolution of the sum is the sum of the
convolutions, so adding the overlapping output pieces reconstructs the exact
linear convolution of the whole stream.

\begin{figure}[H]
\centering
\begin{tikzpicture}[scale=1.0]
  % timeline
  \draw[->,thick] (0,0) -- (11,0) node[right]{$n$};
  % three input blocks of length L=3 units, spaced by L
  \foreach \r/\xo/\col in {0/0/sigblue, 1/3/targetgreen, 2/6/quotientorange} {
    \draw[\col,very thick] (\xo,2.2) -- (\xo+3,2.2);
    \node[\col] at (\xo+1.5,2.5) {block \r};
  }
  \node[left] at (0,2.2) {input};
  % output pieces of length L+N-1 = 3+2 = 5, starting at same offsets, drawn lower, overlapping by N-1=2
  \foreach \r/\xo/\col in {0/0/sigblue, 1/3/targetgreen, 2/6/quotientorange} {
    \draw[\col,very thick,->] (\xo,1.1) -- (\xo+5,1.1);
  }
  \node[left] at (0,1.1) {conv out};
  % overlap brackets
  \draw[trapred,decorate,decoration={brace,amplitude=4pt}] (3,0.85) -- (5,0.85);
  \draw[trapred,decorate,decoration={brace,amplitude=4pt}] (6,0.85) -- (8,0.85);
  \node[trapred] at (4,0.5) {\small add};
  \node[trapred] at (7,0.5) {\small add};
\end{tikzpicture}
\caption{Overlap-add with block length $L=3$ and filter length $N=3$ (so each
convolved block has length $L+N-1=5$). Consecutive output pieces overlap by
$N-1=2$ samples; summing the overlaps reconstructs the exact linear
convolution.}
\label{ch17:fig-ola}
\end{figure}

The recipe, per block:
\begin{enumerate}
  \item Take the next $L$ input samples; zero-pad to length $M \ge L+N-1$ (use a
        power of two for the FFT).
  \item Multiply its FFT by $H = \FT\{h\}$ (computed \emph{once}, ahead of time).
  \item Inverse-FFT to get $L+N-1$ output samples.
  \item Add the first $N-1$ of them to the saved tail from the previous block;
        emit; save this block's last $N-1$ samples as the new tail.
\end{enumerate}

\begin{lstlisting}[language=Python]
import numpy as np

def overlap_add(x, h, L):
    N = len(h)
    M = 1 << int(np.ceil(np.log2(L + N - 1)))   # FFT length, power of 2
    H = np.fft.rfft(h, M)                         # precompute once
    y = np.zeros(len(x) + N - 1)
    for i in range(0, len(x), L):
        blk = x[i:i+L]
        Y = np.fft.irfft(np.fft.rfft(blk, M) * H, M)  # = linear conv of blk*h
        y[i:i+len(Y)] += Y                            # overlap and ADD
    return y[:len(x) + N - 1]
\end{lstlisting}

\subsection{Overlap-save}
\label{ch17:ols}

Overlap-save attacks the same problem from the other side. Instead of zero-padding
the input and adding tails, it deliberately \emph{allows} the circular
wrap-around to corrupt part of each block, then simply throws the corrupted part
away.

You feed the FFT blocks of length $M$ that \emph{overlap the input} by $N-1$
samples (each block reuses the last $N-1$ samples of the previous one). After
circular convolution at length $M$, the first $N-1$ output samples are garbage
(wrapped) and the remaining $M-N+1$ are correct linear-convolution outputs. Save
the good ones, discard the rest. No addition of tails is needed --- hence the
older name ``overlap-discard.''

\begin{table}[H]
\centering
\begin{tabular}{@{}lll@{}}
\toprule
 & \textbf{Overlap-add} & \textbf{Overlap-save} \\
\midrule
Input blocks      & contiguous, no overlap        & overlap by $N-1$ \\
Per block         & zero-pad input                & no zero-pad of input \\
Combine step      & add overlapping output tails  & discard first $N-1$ outputs \\
Bookkeeping       & save and add a tail buffer    & save last $N-1$ input samples \\
Arithmetic        & identical FFT count           & identical FFT count \\
\bottomrule
\end{tabular}
\caption{Overlap-add versus overlap-save. They compute the same answer with the
same number of FFTs; they differ only in how they stitch blocks together. Pick
whichever bookkeeping you find less error-prone.}
\label{ch17:tab-ola-ols}
\end{table}

\begin{trapbox}
The single most common fast-convolution bug is choosing $M$ too small. If
$M < L + N - 1$ the wrap-around eats into real output and you get a quiet,
hard-to-spot distortion --- often a faint echo or a click at block boundaries.
Always assert $M \ge L + N - 1$ in code. The second most common bug is
forgetting that $H = \FT\{h\}$ must be computed at the \emph{same} FFT length
$M$ as the blocks; reuse a precomputed $H$, but recompute it if you change $M$.
\end{trapbox}

\begin{whybox}
\textbf{When is fast convolution actually worth it?} Direct FIR costs about $N$
real MACs per output sample. Overlap-add costs roughly two length-$M$ FFTs and
$M$ complex multiplies per block of $L \approx M$ outputs, i.e.\ on the order of
$\,c\log_2 M\,$ operations per output sample for a small constant $c$. The
crossover is around $N \approx 60$--$100$ taps on a typical CPU: below that,
direct convolution wins (and has lower latency and zero block delay); above it,
FFT convolution pulls steadily ahead, and for $N$ in the thousands it is
dramatically faster. For the wrist tracker's short smoothing filters ($N$ of a
handful to a few dozen), \emph{direct} convolution is both faster and lower
latency --- save fast convolution for long audio impulse responses.
\end{whybox}

\section{Implementing the filter: direct forms}
\label{ch17:directforms}

Set the FFT aside and look at how a filter's difference equation actually becomes
a data-flow of delays, multiplies, and adds. The general LTI difference equation
from Chapter~15 is

\begin{equation}
  y[n] = \sum_{i=0}^{N} b_i\, x[n-i] \;-\; \sum_{k=1}^{M} a_k\, y[n-k],
  \label{ch17:diffeq}
\end{equation}

with feedforward (numerator) coefficients $b_i$ and feedback (denominator)
coefficients $a_k$. An FIR filter is the special case with no feedback
($a_k = 0$); an IIR filter has feedback. There are several ways to wire this up,
all algebraically equal but \emph{numerically different}.

\subsection{Direct Form I and II}

\textbf{Direct Form I} implements Equation~\eqref{ch17:diffeq} literally: one
chain of delays for the past inputs, another for the past outputs, multiply by
the $b$'s and $a$'s, sum. It needs $N + M$ delay registers.

\textbf{Direct Form II} notices that the two delay chains carry redundant
information and merges them, halving the delay storage to $\max(N,M)$. It
computes an intermediate ``state'' signal $w[n]$:
\begin{align}
  w[n] &= x[n] - \sum_{k=1}^{M} a_k\, w[n-k], \\
  y[n] &= \sum_{i=0}^{N} b_i\, w[n-i].
\end{align}

\begin{figure}[H]
\centering
\begin{tikzpicture}[>=Stealth, every node/.style={font=\small}]
  \node (x) {$x[n]$};
  \node[draw,circle,inner sep=1pt,right=10mm of x] (sum1) {$+$};
  \node[draw,minimum width=8mm,minimum height=6mm,right=14mm of sum1] (b0) {$b_0$};
  \node[draw,circle,inner sep=1pt,right=12mm of b0] (sum2) {$+$};
  \node[right=12mm of sum2] (y) {$y[n]$};
  \draw[->] (x) -- (sum1);
  \draw[->] (sum1) -- (b0);
  \draw[->] (b0) -- (sum2);
  \draw[->] (sum2) -- (y);
  % delay chain for w
  \node[draw,minimum width=8mm,minimum height=6mm,below=10mm of b0] (z1) {$z^{-1}$};
  \node[draw,minimum width=8mm,minimum height=6mm,below=10mm of z1] (z2) {$z^{-1}$};
  \draw[->] ($(b0.south)$) -- ++(0,-2mm) -| (z1.north);
  \draw[->] (z1) -- (z2);
  % feedback a1, a2 into sum1
  \node[draw,minimum width=8mm,minimum height=5mm,left=12mm of z1] (a1) {$-a_1$};
  \node[draw,minimum width=8mm,minimum height=5mm,left=12mm of z2] (a2) {$-a_2$};
  \draw[->] (z1.west) -- (a1.east);
  \draw[->] (z2.west) -- (a2.east);
  \draw[->] (a1.west) -| (sum1.south);
  \draw[->] (a2.west) -| ($(sum1.south)+(-3mm,0)$);
  % feedforward b1,b2 into sum2
  \node[draw,minimum width=7mm,minimum height=5mm,right=10mm of z1] (bb1) {$b_1$};
  \node[draw,minimum width=7mm,minimum height=5mm,right=10mm of z2] (bb2) {$b_2$};
  \draw[->] (z1.east) -- (bb1.west);
  \draw[->] (z2.east) -- (bb2.west);
  \draw[->] (bb1.east) -| (sum2.south);
  \draw[->] (bb2.east) -| ($(sum2.south)+(3mm,0)$);
\end{tikzpicture}
\caption{Direct Form II of a second-order section. A single delay line
($z^{-1}$ blocks) holds the shared state $w[n]$; the feedback taps $-a_k$ enter
the left adder and the feedforward taps $b_i$ enter the right adder. Half the
delays of Direct Form I.}
\label{ch17:fig-df2}
\end{figure}

\begin{trapbox}
Direct Form II saves memory but is \emph{worse} numerically for high-order IIR
filters: the shared state $w[n]$ can have a much larger dynamic range than the
input or output, so it overflows (in fixed point) or loses precision (in
floating point) more readily. The practical rule, below, sidesteps this entirely.
\end{trapbox}

\subsection{Cascade of biquads: the practical default}
\label{ch17:biquad}

Never implement a high-order IIR filter as one big Direct Form. The reason is
the same one that makes a high-degree polynomial wretched to evaluate from its
coefficients: the roots (poles and zeros) become wildly sensitive to tiny
coefficient errors, and a filter that was stable on paper can ring or blow up
once its coefficients are rounded to 16 or 24 bits.

The cure is to factor the transfer function into a \emph{cascade} of
second-order sections (``biquads''):

\begin{equation}
  H(z) = g \prod_{s=1}^{S} \frac{b_{0,s} + b_{1,s}\,z^{-1} + b_{2,s}\,z^{-2}}
                                 {1 + a_{1,s}\,z^{-1} + a_{2,s}\,z^{-2}}.
  \label{ch17:cascade}
\end{equation}

Each second-order section holds just one pole pair and one zero pair, so its
coefficients localize the action and stay well-conditioned. You run the sections
one after another, the output of each feeding the next.

\begin{whybox}
Why second order and not first? Real filter poles and zeros come in
\emph{complex-conjugate pairs} (because $h[n]$ is real). A complex pole at
$re^{\jj\theta}$ and its conjugate $re^{-\jj\theta}$ combine into one real
second-order section with $a_1 = -2r\cos\theta$, $a_2 = r^2$. Keeping each
conjugate pair together in its own section is what keeps the coefficients real
and the arithmetic small. This is also why ``biquad'' (bi-quadratic: quadratic
over quadratic) is the universal currency of audio EQ, sensor filters, and DSP
chips --- every parametric EQ band you have ever turned is one biquad.
\end{whybox}

\begin{figure}[H]
\centering
\begin{tikzpicture}[scale=2.6]
  % unit circle and axes
  \draw[gray] (0,0) circle (1);
  \draw[->,thick] (-1.25,0) -- (1.25,0) node[right]{$\mathrm{Re}$};
  \draw[->,thick] (0,-1.25) -- (0,1.25) node[above]{$\mathrm{Im}$};
  % conjugate pole pair at r=0.85, theta=+/-0.9 rad
  \def\rr{0.85}
  \node[specred] at ({\rr*cos(51.6)},{\rr*sin(51.6)}) {$\times$};
  \node[specred] at ({\rr*cos(-51.6)},{\rr*sin(-51.6)}) {$\times$};
  % zero pair on the circle (e.g. at +/-2.4 rad)
  \node[sigblue] at ({cos(137.5)},{sin(137.5)}) {$\circ$};
  \node[sigblue] at ({cos(-137.5)},{sin(-137.5)}) {$\circ$};
  % radius annotation to a pole
  \draw[densely dashed,specred] (0,0) -- ({\rr*cos(51.6)},{\rr*sin(51.6)});
  \node[specred,font=\scriptsize] at (0.32,0.28) {$r$};
  \node[specred,font=\scriptsize] at (0.62,0.78) {pole $re^{\jj\theta}$};
\end{tikzpicture}
\caption{Pole-zero plot of one biquad. The complex-conjugate pole pair (red
$\times$) at $re^{\pm\jj\theta}$ gives the real denominator
$1+a_1 z^{-1}+a_2 z^{-2}$ with $a_1=-2r\cos\theta$, $a_2=r^2$; the conjugate zero
pair (blue $\circ$) gives the numerator. Poles inside the unit circle keep the
section stable; the closer they sit to the circle, the sharper and more
fixed-point-sensitive the resonance.}
\label{ch17:fig-polezero}
\end{figure}

\begin{lstlisting}[language=C]
// One biquad, Transposed Direct Form II (good numerics, 2 state vars).
typedef struct { float b0,b1,b2,a1,a2; float z1,z2; } biquad;

float biquad_step(biquad *q, float x) {
    float y = q->b0 * x + q->z1;
    q->z1   = q->b1 * x - q->a1 * y + q->z2;
    q->z2   = q->b2 * x - q->a2 * y;
    return y;
}
// A high-order filter is just an array of these run in sequence.
float cascade_step(biquad *qs, int S, float x) {
    for (int s = 0; s < S; s++) x = biquad_step(&qs[s], x);
    return x;
}
\end{lstlisting}

The \emph{transposed} Direct Form II shown here is the workhorse: it has the
same low storage as Direct Form II but much better numerical behavior, because
its two state variables stay closer to the signal's natural range. It is what
you will find inside almost every audio library's biquad.

\section{Fixed-point and numerical issues}
\label{ch17:fixedpoint}

On a desktop you filter in 32- or 64-bit floating point and rarely think about
arithmetic error. On the small microcontroller that will live on the wrist
strap, you may have no floating-point unit at all, or a slow one. Then you filter
in \emph{fixed point}: integers interpreted as fractions, with an implied binary
point. Understanding what goes wrong there is the difference between a filter
that works on the bench and one that hisses, clips, or drifts in the field.

\subsection{What fixed point is}

A fixed-point format $Qm.f$ stores a number as an integer scaled by $2^{-f}$:
the integer $X$ represents the real value $X \cdot 2^{-f}$, with $m$ integer bits
and $f$ fractional bits. For example $Q1.15$ uses a 16-bit signed integer to
represent values in $[-1, 1)$ with a resolution of $2^{-15} \approx 3\times
10^{-5}$. Adding two $Q1.15$ numbers is just integer addition; \emph{multiplying}
two of them gives a $Q2.30$ result that you must shift back down by 15 bits to
return to $Q1.15$.

\begin{connectionbox}
This is exactly the quantization story of Chapter~19, applied not to the input
sample but to every internal coefficient and intermediate result. Each rounding
injects a tiny error that, statistically, behaves like added noise (Chapter~20):
uniform on $\pm \tfrac12$ LSB, with variance $\Delta^2/12$ where $\Delta=2^{-f}$
is the quantization step. A filter is a chain of such injections, and the
question is whether they accumulate or wash out.
\end{connectionbox}

\subsection{The three things that bite}

\paragraph{1. Coefficient quantization.} Rounding the $b_i$, $a_k$ to fixed
point moves the poles and zeros. For a sharp, high-order filter the poles sit
close together and close to the unit circle, where small coefficient nudges
cause big pole movements --- exactly the ill-conditioning that the biquad cascade
fixes. A direct-form 8th-order filter quantized to 16 bits can have its poles
shoved outside the unit circle and \emph{go unstable}; the same filter as four
biquads stays put.

\paragraph{2. Roundoff (product) noise.} Every multiply produces more bits than
fit, so the result is rounded, adding noise. In an IIR filter this noise
circulates through the feedback and can be amplified near resonant peaks. The
defense is to keep enough fractional bits and to scale section gains so signals
use the available range without overflowing.

\paragraph{3. Overflow and limit cycles.} If an intermediate sum exceeds the
representable range it \emph{wraps} (in two's-complement) from the largest
positive value to the largest negative --- a catastrophic, full-scale glitch.
Worse, the feedback in an IIR filter can sustain a small self-perpetuating
oscillation called a \emph{limit cycle}: with a constant or zero input, the
rounded output settles into a tiny periodic wiggle instead of going to rest.

\begin{figure}[H]
\centering
\begin{tikzpicture}
\begin{axis}[
  width=11cm, height=5cm,
  xlabel={sample $n$}, ylabel={output},
  domain=0:40, samples=41,
  ymin=-1.2, ymax=1.2, xmin=0, xmax=40,
  axis lines=left, enlarge x limits=false,
  legend pos=north east, legend style={font=\footnotesize},
]
% ideal: decaying ringing to zero
\addplot[sigblue, very thick] {0.95^x * cos(deg(0.6*x))};
\addlegendentry{ideal (float): decays to 0}
% quantized: decays then locks into a small limit cycle
\addplot[specred, very thick, dashed] {max(abs(0.95^x*cos(deg(0.6*x))), 0.06)*sign(cos(deg(0.6*x)))};
\addlegendentry{fixed point: limit cycle}
\addplot[gray, dotted, thick] coordinates {(0,0.06)(40,0.06)};
\addplot[gray, dotted, thick] coordinates {(0,-0.06)(40,-0.06)};
\end{axis}
\end{tikzpicture}
\caption{A limit cycle in an IIR filter. In floating point the resonator's
ringing decays smoothly to zero (solid). In fixed point, rounding can trap the
output in a small sustained oscillation (dashed) that never dies --- audible as
a faint idle tone, or as a sensor reading that refuses to settle.}
\label{ch17:fig-limitcycle}
\end{figure}

\begin{focusbox}
Practical fixed-point survival kit:
\begin{itemize}
  \item Implement IIR filters as cascaded \emph{biquads}, never one big direct
        form.
  \item Use \emph{transposed Direct Form II} sections.
  \item Accumulate products in a wider register (e.g.\ 32- or 64-bit accumulator
        for 16-bit data) and round only once, at the end.
  \item \emph{Scale} each section so internal signals fill the range without
        overflowing; saturate (clamp) rather than wrap on overflow.
  \item For FIR filters, prefer fixed point freely --- with no feedback they
        cannot have limit cycles or go unstable, only add bounded roundoff noise.
\end{itemize}
\end{focusbox}

\begin{trapbox}
\emph{Saturating} arithmetic (clamp to the max/min on overflow) turns a
catastrophic wrap-around glitch into a graceful, soft distortion. Most DSP
instruction sets have saturating add/multiply for exactly this reason. If your
platform wraps by default, add an explicit clamp on the accumulator before
storing the output. One missing clamp is the classic cause of an occasional loud
``pop'' in an otherwise clean audio path.
\end{trapbox}

\section{Latency: the clock you cannot cheat}
\label{ch17:latency}

For a musical controller, latency is the single most important number in the
whole system. A performer feels the gap between moving and hearing the result in
their fingertips, and beyond a threshold the instrument stops feeling like part
of the body. So let us account for every microsecond a filter costs.

\subsection{Group delay: how late the filter makes the signal}

A filter's frequency response $H(e^{\jj\omega}) = \abs{H}\,e^{\jj\phi(\omega)}$
has a phase $\phi(\omega)$. The \emph{group delay} is the negative slope of that
phase,
\begin{equation}
  \tau_g(\omega) = -\frac{\dd \phi(\omega)}{\dd\omega},
  \label{ch17:groupdelay}
\end{equation}
measured in samples. It is, intuitively, how many samples late the envelope of a
narrow band of frequencies near $\omega$ comes out of the filter. (We met this in
Chapter~14; here it becomes a budget line.)

\begin{happeningbox}
For a \emph{symmetric, linear-phase FIR} of length $N$, the phase is a straight
line and the group delay is the \emph{same at every frequency}:
\[
  \tau_g = \frac{N-1}{2}\ \text{samples}.
\]
Every frequency is delayed equally, so the waveform shape is preserved --- but it
is preserved \emph{late}, by half the filter length. A 101-tap linear-phase
low-pass at $f_s = 200$~Hz delays everything by $50$ samples, i.e.\ $50/200 =
0.25$~s --- a quarter second! That is unplayable. Linear phase is gorgeous for
offline analysis and ruinous for live control.
\end{happeningbox}

\begin{intuitionbox}
Linear-phase FIR is symmetric, so its ``center of mass'' is in the middle. To
emit the centered output you must wait until the middle tap lines up with the
newest sample --- which means waiting $(N-1)/2$ samples after the event you care
about. The sharper the filter (the more taps), the longer the wait. There is no
trick around it: \textbf{a sharp linear-phase filter is inherently laggy.}
\end{intuitionbox}

\subsection{IIR buys low latency with phase distortion}

An IIR filter reaches a given sharpness with far fewer coefficients than an FIR,
and its group delay can be a handful of samples even for a steep response. The
price is \emph{nonlinear phase}: different frequencies are delayed by different
amounts, so transients smear in shape. For control signals (we mostly care
\emph{when} a gesture happens, not the pristine shape of the waveform) this trade
is usually worth it: an IIR low-pass with $\tau_g$ of a few samples beats a
linear-phase FIR with $\tau_g$ of fifty.

\begin{figure}[H]
\centering
\begin{tikzpicture}
\begin{axis}[
  width=12cm, height=5.2cm,
  xlabel={normalized frequency $\omega$ (rad/sample)}, ylabel={group delay (samples)},
  domain=0:3.1416, samples=120,
  ymin=0, ymax=60, xmin=0, xmax=3.1416,
  axis lines=left, enlarge x limits=false,
  legend pos=north west, legend style={font=\footnotesize},
  xtick={0,1.5708,3.1416}, xticklabels={$0$,$\pi/2$,$\pi$},
]
% linear-phase FIR: constant 50 samples
\addplot[sigblue, very thick] {50};
\addlegendentry{linear-phase FIR (101 taps)}
% IIR: small, frequency-dependent, peaking near cutoff ~ 1.0
\addplot[specred, very thick, dashed] {3 + 7*exp(-((x-1.0)*2.2)^2)};
\addlegendentry{4th-order IIR (biquads)}
\end{axis}
\end{tikzpicture}
\caption{Group delay versus frequency for two filters of comparable sharpness.
The FIR delays every frequency by a flat $50$ samples (linear phase). The IIR
delays by only a few samples but unevenly, peaking near its cutoff (nonlinear
phase). For live control the IIR's low, if uneven, delay usually wins.}
\label{ch17:fig-groupdelay}
\end{figure}

\subsection{The full latency budget}

End-to-end latency for the wrist controller is a sum, and the filter is only one
term:
\begin{equation}
  \tau_{\text{total}} =
    \underbrace{\tau_{\text{sensor}}}_{\text{ADC + internal LPF}}
    + \underbrace{\tau_{\text{block}}}_{L\,T_s\ \text{if buffered}}
    + \underbrace{\tau_{\text{filter}}}_{\tau_g \cdot T_s}
    + \underbrace{\tau_{\text{compute}}}_{\text{DSP time}}
    + \underbrace{\tau_{\text{link}}}_{\text{MIDI/BLE}}
    + \underbrace{\tau_{\text{synth}}}_{\text{sound engine}}.
  \label{ch17:budget}
\end{equation}

\begin{goalbox}
Musicians reliably notice action-to-sound latency above roughly $10$~ms, and
$20$--$30$~ms feels sluggish for percussive playing. Budget accordingly:
\begin{itemize}
  \item \textbf{Sensor LPF:} MEMS IMUs have a built-in anti-alias/low-pass with
        its own group delay --- read the datasheet; a low bandwidth setting can
        cost several ms. Choose a bandwidth that just covers your gesture
        frequencies (motion energy is mostly below $20$~Hz).
  \item \textbf{Blocking:} go sample-by-sample on the control path; do not buffer
        the gesture detector.
  \item \textbf{Filter:} use short IIR smoothers ($\tau_g$ of a few samples).
        Avoid long linear-phase FIRs on the live path. At $f_s = 200$~Hz, each
        sample is $5$~ms, so even a 5-sample group delay is $25$~ms --- so
        \emph{sample faster} ($f_s \ge 500$~Hz) on the control path to shrink the
        per-sample time.
  \item \textbf{Link:} classic serial MIDI adds about $1$~ms per message;
        Bluetooth LE MIDI can add $10$--$30$~ms or more and is jittery ---
        prefer USB or wired for serious playing.
\end{itemize}
The takeaway: filtering latency is one of several terms, and the wireless
link often dwarfs it. Design the whole chain, not just the filter.
\end{goalbox}

\begin{historybox}
The $10$~ms figure is not folklore. Studies of musicians (e.g.\ work by
Wessel and Wright at Berkeley's CNMAT, and later latency-perception studies)
consistently find that performers detect and are disturbed by action-to-sound
delays in the $10$~ms range, with skilled percussionists more sensitive still.
It is the same order as the acoustic delay of standing a couple of meters from
your own amplifier --- a delay real musicians already learn to hate.
\end{historybox}

\subsection{Zero-phase filtering --- and why you cannot use it live}

There is a beloved offline trick: filter the signal forward, then filter the
result \emph{backward} with the same filter (\texttt{filtfilt} in SciPy/MATLAB).
The backward pass cancels the forward pass's phase exactly, giving \emph{zero}
net phase distortion and zero net delay --- and it squares the magnitude
response, sharpening it.

\begin{trapbox}
Zero-phase filtering requires running the filter backward in time, which means
you must have the \emph{entire signal} already. It is fundamentally non-causal
and \textbf{cannot be used in real time}. It is perfect for cleaning a stored
IMU log offline, useless on the live wrist stream. Do not let a clean offline
result fool you into thinking the same filter will work live --- the live version
will have all the delay back, plus phase distortion.
\end{trapbox}

\section{Moving average: the workhorse smoother}
\label{ch17:movingavg}

Almost every real sensor pipeline starts with a moving average, because it is the
cheapest possible low-pass filter and it is exactly the right tool for knocking
down broadband noise on a slowly-varying signal. The length-$L$ moving average is

\begin{equation}
  y[n] = \frac{1}{L}\sum_{k=0}^{L-1} x[n-k].
  \label{ch17:ma}
\end{equation}

It is an FIR filter whose $L$ taps are all $1/L$ --- a rectangular impulse
response. We computed its frequency response in Chapter~14; its magnitude is the
\emph{aliased sinc} (Dirichlet kernel):
\begin{equation}
  \abs{H(e^{\jj\omega})}
  = \frac{1}{L}\left|\frac{\sin(L\omega/2)}{\sin(\omega/2)}\right|.
  \label{ch17:ma-mag}
\end{equation}

\begin{figure}[H]
\centering
\begin{tikzpicture}
\begin{axis}[
  width=12cm, height=5.2cm,
  xlabel={normalized frequency $\omega$ (rad/sample)}, ylabel={$|H|$},
  domain=0.001:3.1416, samples=180,
  ymin=0, ymax=1.05, xmin=0, xmax=3.1416,
  axis lines=left, enlarge x limits=false,
  legend pos=north east, legend style={font=\footnotesize},
  xtick={0,1.5708,3.1416}, xticklabels={$0$,$\pi/2$,$\pi$},
]
\addplot[sigblue, very thick] {abs(sin(deg(4*x/2))/(4*sin(deg(x/2))))};
\addlegendentry{$L=4$}
\addplot[specred, very thick, dashed] {abs(sin(deg(8*x/2))/(8*sin(deg(x/2))))};
\addlegendentry{$L=8$}
\addplot[targetgreen, very thick, dotted] {abs(sin(deg(16*x/2))/(16*sin(deg(x/2))))};
\addlegendentry{$L=16$}
\end{axis}
\end{tikzpicture}
\caption{Magnitude response of the length-$L$ moving average. Longer windows
push the first null lower (more smoothing) but the stopband lobes stay high
($\approx -13$~dB) --- the moving average is a gentle low-pass with leaky
stopband, not a sharp filter. The nulls sit at $\omega = 2\pi m/L$.}
\label{ch17:fig-ma}
\end{figure}

\begin{happeningbox}
The moving average is a \emph{terrible} frequency-selective filter --- its first
sidelobe is only $13$~dB down, so it lets plenty of high frequency leak through.
But it is a \emph{wonderful} smoother of broadband noise, because averaging $L$
independent noisy samples reduces the noise standard deviation by $\sqrt{L}$
(Chapter~20) while passing slow signal almost untouched. When your goal is
``make this jittery sensor reading calmer,'' not ``isolate a precise frequency
band,'' the moving average is usually the right and cheapest answer.
\end{happeningbox}

\subsection{The running-sum trick: \texorpdfstring{$O(1)$}{O(1)} per sample}

The naive moving average does $L$ adds per output. But a moving \emph{sum}
changes by only two terms each step: add the new sample, subtract the one that
just fell off the window. So you can compute it in \emph{constant} time
regardless of $L$:

\begin{equation}
  S[n] = S[n-1] + x[n] - x[n-L], \qquad y[n] = S[n]/L.
  \label{ch17:runsum}
\end{equation}

\begin{lstlisting}[language=C]
// O(1) moving average via running sum. ring[] holds the last L samples.
typedef struct { float ring[64]; int L, pos; float sum; } movavg;

float movavg_step(movavg *m, float x) {
    m->sum += x - m->ring[m->pos];  // add new, drop oldest
    m->ring[m->pos] = x;
    m->pos = (m->pos + 1) % m->L;
    return m->sum / m->L;
}
\end{lstlisting}

\begin{trapbox}
The running-sum recurrence is numerically fragile in fixed point or with a
constant DC offset: rounding errors and bias accumulate in $S$ forever, since the
two correction terms never quite cancel in finite precision. Over millions of
samples $S$ can drift. Use a wide integer accumulator (or periodically recompute
the sum from scratch), and prefer it only where the per-sample savings matter.
For short $L$ (a handful of taps) just sum directly --- it is robust and the cost
is trivial.
\end{trapbox}

\subsection{The exponential moving average (one-pole IIR)}

The moving average has a hard, finite window. Its cousin the \emph{exponential}
moving average (EMA) weights the past with a smooth exponential decay and is the
simplest IIR filter there is:

\begin{equation}
  y[n] = \alpha\, x[n] + (1-\alpha)\, y[n-1], \qquad 0 < \alpha < 1.
  \label{ch17:ema}
\end{equation}

One multiply, one add, one state variable. Smaller $\alpha$ means heavier
smoothing (a longer effective memory). Its single pole sits at $z = 1-\alpha$,
its time constant is roughly $\tau \approx 1/\alpha$ samples, and its $-3$~dB
cutoff is approximately $\omega_c \approx \alpha$ for small $\alpha$.

\begin{figure}[H]
\centering
\begin{tikzpicture}
\begin{axis}[
  width=12cm, height=5cm,
  xlabel={sample $n$}, ylabel={step response $y[n]$},
  domain=0:40, samples=41,
  ymin=0, ymax=1.1, xmin=0, xmax=40,
  axis lines=left, enlarge x limits=false,
  legend pos=south east, legend style={font=\footnotesize},
]
\addplot[sigblue, very thick] {1 - (1-0.3)^x};
\addlegendentry{$\alpha=0.3$ (light smoothing)}
\addplot[specred, very thick, dashed] {1 - (1-0.1)^x};
\addlegendentry{$\alpha=0.1$ (heavy smoothing)}
\addplot[gray, dotted, thick] coordinates {(0,0.95)(40,0.95)};
\end{axis}
\end{tikzpicture}
\caption{Step response of the exponential moving average for two values of
$\alpha$. The output rises toward the new level as $1-(1-\alpha)^n$; smaller
$\alpha$ smooths harder but lags more, taking longer to reach the $95\%$ line
(dotted). The effective memory is about $1/\alpha$ samples.}
\label{ch17:fig-ema}
\end{figure}

\begin{focusbox}
\textbf{Moving average vs.\ EMA --- which to reach for.}
\begin{itemize}
  \item \emph{Moving average:} flat weighting, sharper transition for the same
        memory, exact known group delay $(L-1)/2$, but needs an $L$-sample
        buffer. Good when you want a defined window and linear phase.
  \item \emph{EMA:} one state variable, no buffer, trivially cheap, instantly
        tunable by one knob $\alpha$. Group delay is roughly $(1-\alpha)/\alpha$
        samples. Good when memory is tight and you want a smooth, gently-decaying
        response --- the default for embedded sensor smoothing.
\end{itemize}
\end{focusbox}

\section{Median filter: killing spikes without smearing edges}
\label{ch17:median}

Linear filters --- moving average, EMA, any FIR/IIR --- share a fatal weakness
against \emph{impulsive} noise. A single wild sample (a sensor glitch, an
electrical spike, a dropped I2C read returning garbage) gets \emph{spread} by a
linear filter: the moving average smears one bad value across $L$ outputs instead
of removing it. Worse, it pulls the average toward the outlier.

The \emph{median} filter replaces each sample with the median of a sliding window
of $L$ samples (odd $L$):
\begin{equation}
  y[n] = \operatorname{median}\bigl(x[n-k],\ \dots,\ x[n+k]\bigr),
  \qquad L = 2k+1.
  \label{ch17:median-def}
\end{equation}

The median is a \emph{nonlinear} operation --- it is not a convolution, has no
frequency response, and does not obey superposition. But it has a magical
property: it is immune to outliers. A single spike, no matter how huge, cannot be
the median of a window in which it is a minority, so it is simply discarded ---
while genuine edges and steps in the signal are preserved sharply, not blurred.

\begin{figure}[H]
\centering
\begin{tikzpicture}
\begin{axis}[
  width=12cm, height=5.5cm,
  xlabel={sample $n$}, ylabel={value},
  xmin=0, xmax=19, ymin=-0.5, ymax=4.5,
  axis lines=left, enlarge x limits=false,
  legend pos=north west, legend style={font=\footnotesize},
]
% raw: a step from 1 to 3 at n=10, with two spikes
\addplot[gray, thick, mark=*, mark size=1pt] coordinates {
 (0,1)(1,1)(2,1)(3,4)(4,1)(5,1)(6,1)(7,1)(8,1)(9,1)(10,3)(11,3)(12,3)(13,3)(14,0)(15,3)(16,3)(17,3)(18,3)(19,3)};
\addlegendentry{raw (step + spikes)}
% median L=3 result: spikes removed, edge preserved
\addplot[specred, very thick] coordinates {
 (0,1)(1,1)(2,1)(3,1)(4,1)(5,1)(6,1)(7,1)(8,1)(9,1)(10,3)(11,3)(12,3)(13,3)(14,3)(15,3)(16,3)(17,3)(18,3)(19,3)};
\addlegendentry{median, $L=3$}
% moving avg L=3: spikes smeared, edge blurred
\addplot[sigblue, very thick, dashed] coordinates {
 (0,1)(1,1)(2,2)(3,2)(4,2)(5,1)(6,1)(7,1)(8,1)(9,1.67)(10,2.33)(11,3)(12,3)(13,2)(14,2)(15,2)(16,3)(17,3)(18,3)(19,3)};
\addlegendentry{moving avg, $L=3$}
\end{axis}
\end{tikzpicture}
\caption{Median versus moving average on a step contaminated by two spikes (at
$n=3$ and $n=14$). The median (solid red) removes both spikes completely and
keeps the step edge razor-sharp. The moving average (dashed blue) smears each
spike into its neighbors and rounds off the edge.}
\label{ch17:fig-median}
\end{figure}

\begin{whybox}
Why does the median preserve edges while the mean blurs them? At a step, the
window straddling the edge contains some ``low'' samples and some ``high''
samples. The mean returns an intermediate value --- a ramp across the edge. The
median returns whichever side is the majority --- it stays low until the window
is more than half past the edge, then jumps. The output is a clean step, just
shifted by the window center. Same reason it kills spikes: a spike is a minority
of one in its window, so the median never selects it.
\end{whybox}

\begin{lstlisting}[language=Python]
import numpy as np

def median_filter(x, L):       # L odd
    k = L // 2
    xp = np.pad(x, k, mode='edge')
    return np.array([np.median(xp[i:i+L]) for i in range(len(x))])
\end{lstlisting}

\begin{trapbox}
The median filter is more expensive (sorting a window is $O(L\log L)$ naively,
or $O(L)$ with a maintained sorted structure) and \emph{nonlinear}: cascading two
median filters is not the same as one, and you cannot reason about it with a
frequency response. It also rounds off thin peaks narrower than $k$ samples ---
if a genuine narrow feature is real signal, the median will eat it just like a
spike. Use it specifically against impulsive noise, not as a general smoother.
\end{trapbox}

\begin{applicationbox}
Median filtering is everywhere outliers appear: removing salt-and-pepper noise
from images (a 2D median, Chapter~32), de-glitching GPS fixes, cleaning ECG
baselines of motion artifacts, and --- for us --- rejecting the occasional
corrupt sample from a flaky I2C/SPI sensor read. A common, cheap, and very
effective combination is a short \emph{median} (length 3 or 5) to kill spikes,
followed by a moving average or EMA to smooth the residual broadband jitter.
\end{applicationbox}

\section{Putting a real sensor pipeline together}
\label{ch17:pipeline}

We now have all the pieces. A robust, low-latency de-noising front end for a
streaming sensor is almost always the same small stack, in this order:

\begin{enumerate}
  \item \textbf{De-glitch} with a short median (length 3 or 5) to discard
        impulsive spikes and dropped-read garbage \emph{before} they can
        contaminate any averaging.
  \item \textbf{Smooth} with an EMA or short moving average to knock down the
        broadband jitter that remains, choosing the cutoff to just cover your
        gesture frequencies (motion energy lives below $\sim 20$~Hz).
  \item \textbf{Remove drift / DC} with a gentle high-pass (a one-pole DC
        blocker $y[n] = x[n]-x[n-1]+\rho\,y[n-1]$, $\rho$ near $1$) if you only
        care about \emph{changes}, e.g.\ taps and motion onsets, not absolute
        orientation.
\end{enumerate}

\begin{lstlisting}[language=C]
// Streaming sensor front end: median(3) -> EMA -> DC block.
// State persists across calls; one sample in, one out; minimal latency.
#include <math.h>
typedef struct { float m0,m1,m2; float ema; float xp,yp; } pipe;

static float median3(float a, float b, float c){
    float hi = fmaxf(a,fmaxf(b,c)), lo = fminf(a,fminf(b,c));
    return a + b + c - hi - lo;            // the middle value
}
float pipe_step(pipe *p, float x, float alpha, float rho){
    p->m2 = p->m1; p->m1 = p->m0; p->m0 = x;     // shift median window
    float md = median3(p->m0, p->m1, p->m2);     // 1) de-glitch
    p->ema  = alpha*md + (1.0f-alpha)*p->ema;    // 2) smooth (EMA)
    float dc = p->ema - p->xp + rho*p->yp;       // 3) DC block (high-pass)
    p->xp = p->ema; p->yp = dc;
    return dc;
}
\end{lstlisting}

\begin{goalbox}
This three-stage front end is exactly what runs first on every channel of the
wrist tracker's IMU --- three accelerometer axes and three gyro axes. It is
sample-by-sample (no block latency), costs a handful of operations per sample,
has no large group delay, and is robust to the glitchy reads you \emph{will} get
from a sensor bus on a moving, sweaty wrist. Everything downstream --- orientation
fusion (Chapter~36), gesture features (Chapter~38) --- assumes its input has
already been cleaned by something like this.

And the caveat from the start of the book still holds: this pipeline
makes the IMU's \emph{coarse motion} signal clean and responsive, but no amount
of filtering conjures information the sensor never captured. A wrist IMU can give
you crisp, low-latency detection of \emph{discrete} gestures (a flick, a twist, a
tap) and continuous \emph{wrist} orientation --- it cannot recover individual
finger joint angles. For finger-level control you need EMG/FMG on the forearm or
sensors on the fingers (Chapter~37). Clean signals, and claims that match what
the sensor can actually measure.
\end{goalbox}

\section{Summary}
\label{ch17:summary}

\begin{itemize}
  \item \textbf{Streaming} means bounded memory on an endless signal;
        \textbf{real-time} means meeting a per-sample deadline. Causality forbids
        peeking ahead, so every filter delays.
  \item \textbf{Sample-by-sample} processing minimizes latency and suits IIR and
        smoothers; \textbf{block} processing enables FFT-based work but adds one
        block of delay.
  \item \textbf{Fast convolution} (overlap-add / overlap-save) makes long FIR
        filtering cheap via the FFT, provided $M \ge L+N-1$ so circular
        convolution equals linear. It wins above $\sim 60$--$100$ taps; below
        that, direct convolution is faster and lower latency.
  \item Implement IIR filters as \textbf{cascaded transposed-Direct-Form-II
        biquads}, never one big direct form, to stay stable and accurate ---
        especially in \textbf{fixed point}, where coefficient quantization,
        roundoff noise, overflow, and limit cycles are the dangers.
  \item \textbf{Latency} is a budget. Linear-phase FIR delays by $(N-1)/2$
        samples; short IIR delays by a few; zero-phase \texttt{filtfilt} is
        offline only. For a musical controller, keep total action-to-sound delay
        near $10$~ms.
  \item For real sensor data, the practical de-noising stack is a short
        \textbf{median} to kill spikes, then a \textbf{moving average / EMA} to
        smooth, optionally a \textbf{DC blocker} to remove drift --- cheap,
        robust, and low-latency.
\end{itemize}

\section{Exercises}
\label{ch17:exercises}

\begin{exercise}
A streaming filter runs on a microcontroller at $f_s = 500$~Hz. Each output
sample takes $1.2$~ms of compute. Is this real-time? What is the maximum FIR
length you could afford if each tap costs $2\,\mu$s and you must also leave
$30\%$ of the sample period free for other work?
\end{exercise}

\begin{exercise}
You filter a stream with an FIR of length $N = 200$ using overlap-add with input
blocks of length $L = 256$. (a) What is the minimum FFT length $M$? Round up to a
power of two. (b) How many samples of the overlapping output tail must be saved
and added between consecutive blocks? (c) What is the block-induced latency in
samples, and in milliseconds at $f_s = 48$~kHz?
\end{exercise}

\begin{exercise}
Show by direct expansion that circular convolution of a length-$3$ signal
$x = [x_0,x_1,x_2]$ with a length-$3$ filter $h=[h_0,h_1,h_2]$ at DFT length
$M=4$ does \emph{not} equal their linear convolution, but at $M=5$ it does.
Identify exactly which output sample is corrupted at $M=4$ and why.
\end{exercise}

\begin{exercise}
A length-$101$ symmetric linear-phase low-pass FIR runs at $f_s = 200$~Hz.
(a) What is its group delay in samples and in milliseconds? (b) If a musician
needs total latency under $10$~ms and this filter alone exceeds it, name two
concrete changes that would bring the filter's contribution under $5$~ms, and
state the cost of each.
\end{exercise}

\begin{exercise}
An exponential moving average uses $\alpha = 0.1$. (a) Where is its pole?
(b) Estimate its time constant in samples and its approximate $-3$~dB cutoff in
rad/sample. (c) At $f_s = 200$~Hz, what is that cutoff in Hz? (d) Roughly how many
samples does it take for a step input to reach $95\%$ of its final value?
\end{exercise}

\begin{exercise}
Implement the running-sum moving average of Equation~\eqref{ch17:runsum} in any
language and feed it a constant input of $1.0$ for $10^7$ samples using
\texttt{float32} arithmetic. Compare the final output to the exact value $1.0$.
Explain the discrepancy in terms of accumulated roundoff, and show that a
\texttt{float64} or integer accumulator fixes it.
\end{exercise}

\begin{exercise}
Take the signal in Figure~\ref{ch17:fig-median} (a step from $1$ to $3$ at
$n=10$ with a spike of value $4$ at $n=3$ and a spike of value $0$ at $n=14$). By
hand, compute the output of (a) a length-$3$ median filter and (b) a length-$3$
moving average at samples $n = 3, 4, 13, 14, 15$. Confirm the median removes the
spikes and the average smears them.
\end{exercise}

\begin{exercise}
A second-order IIR resonator has poles at $r e^{\pm \jj\theta}$ with $r = 0.99$,
$\theta = 0.5$. (a) Write the real biquad denominator coefficients $a_1, a_2$.
(b) Explain why this filter is especially prone to fixed-point trouble (limit
cycles, coefficient sensitivity), referencing the pole radius. (c) What single
implementation choice most reduces the risk?
\end{exercise}

\begin{exercise}
Explain in one paragraph each why (a) you cannot use SciPy's \texttt{filtfilt}
zero-phase filtering on the live wrist-controller stream, and (b) a moving average
is a poor choice when you need to isolate a narrow frequency band but an
excellent choice when you simply want to reduce broadband sensor noise.
\end{exercise}

\begin{exercise}
Design (on paper) the de-noising front end for the wrist controller's gyroscope
channel. The gyro is sampled at $f_s = 500$~Hz, occasionally returns a single
garbage sample from a bad SPI read, and carries broadband noise plus a slow
temperature-dependent bias drift. Specify each stage (type, length or $\alpha$),
justify your choices against latency and the gesture bandwidth ($< 20$~Hz), and
estimate the total group delay of your chain in milliseconds.
\end{exercise}

\begin{exercise}
(Harder.) The transposed Direct Form II biquad in the C listing of
Section~\ref{ch17:biquad} uses two state variables $z_1, z_2$. Trace through the
update equations for an impulse input $x = [1,0,0,\dots]$ with coefficients
$b_0=1,\ b_1=b_2=0,\ a_1=-0.9,\ a_2=0$ (a one-pole filter dressed as a biquad).
Confirm that the output is the expected geometric decay $y[n] = 0.9^n$, and state
the pole location.
\end{exercise}
